@article{DrWatson2020,
 pages         = {2673},
 author        = {Datseris, George and Isensee, Jonas and Pech, Sebastian and Gál, Tamás},
 journal       = {Journal of Open Source Software},
 title         = {DrWatson: the perfect sidekick for your scientific inquiries},
 publisher     = {The Open Journal},
 number        = {54},
 doi           = {10.21105/joss.02673},
 year          = {2020},
 url           = {https://doi.org/10.21105/joss.02673},
 volume        = {5}
}

@article{Makie,
  doi = {10.21105/joss.03349},
  url = {https://doi.org/10.21105/joss.03349},
  year = {2021},
  publisher = {The Open Journal},
  volume = {6},
  number = {65},
  pages = {3349},
  author = {Simon Danisch and Julius Krumbiegel},
  title = {{Makie.jl}: Flexible high-performance data visualization for {Julia}},
  journal = {Journal of Open Source Software}
}

@article{DynamicalSystems.jl-2018,
 pages         = {598},
 author        = {Datseris, George},
 month         = {mar},
 journal       = {Journal of Open Source Software},
 title         = {DynamicalSystems.jl: A Julia software library for chaos and nonlinear dynamics},
 number        = {23},
 doi           = {10.21105/joss.00598},
 year          = {2018},
 url           = {https://doi.org/10.21105/joss.00598},
 volume        = {3}
}

@book{DatserisParlitz2022,
 publisher     = {Springer Nature},
 doi           = {10.1007/978-3-030-91032-7},
 language      = {en},
 address       = {Cham, Switzerland},
 author        = {Datseris, George and Parlitz, Ulrich},
 year          = {2022},
 url           = {https://doi.org/10.1007/978-3-030-91032-7},
 title         = {Nonlinear dynamics: A concise introduction interlaced with code}
}

@article{JSSv098i16,
 pages         = {1--30},
 author        = {Besançon, Mathieu and Papamarkou, Theodore and Anthoff, David and Arslan, Alex and Byrne, Simon and Lin, Dahua and Pearson, John},
 journal       = {Journal of Statistical Software},
 title         = {Distributions.jl: Definition and Modeling of Probability Distributions in the JuliaStats Ecosystem},
 number        = {16},
 keywords      = {Julia; distributions; modeling; interface; mixture; KDE; sampling; probabilistic programming; inference},
 doi           = {10.18637/jss.v098.i16},
 year          = {2021},
 url           = {https://www.jstatsoft.org/v098/i16},
 volume        = {98}
}

@misc{Distributions.jl-2019,
 doi           = {10.5281/zenodo.2647458},
 author        = {Lin, Dahua and White, John  Myles and Byrne, Simon and Bates, Douglas and Noack, Andreas and Pearson, John and Arslan, Alex and Squire, Kevin and Anthoff, David and Papamarkou, Theodore and Besançon, Mathieu and Drugowitsch, Jan and Schauer, Moritz and contributors},
 year          = {2019},
 url           = {https://doi.org/10.5281/zenodo.2647458},
 month         = {July},
 title         = {{JuliaStats/Distributions.jl: a Julia package for probability distributions and associated functions}}
}

@article{BenchmarkTools.jl-2016,
 primaryClass  = {cs.PF},
 author        = {{Chen}, Jiahao and {Revels}, Jarrett},
 adsnote       = {Provided by the SAO/NASA Astrophysics Data System},
 month         = {Aug},
 journal       = {arXiv e-prints},
 archivePrefix = {arXiv},
 title         = {{Robust benchmarking in noisy environments}},
 eid           = {arXiv:1608.04295},
 keywords      = {Computer Science - Performance, 68N30, B.8.1, D.2.5},
 eprint        = {1608.04295},
 year          = {2016},
 adsurl        = {https://ui.adsabs.harvard.edu/abs/2016arXiv160804295C}
}

@misc{ma2021modelingtoolkit,
 primaryClass  = {cs.MS},
 author        = {Ma, Yingbo and Gowda, Shashi and Anantharaman, Ranjan and Laughman, Chris and Shah, Viral and Rackauckas, Chris},
 eprint        = {2103.05244},
 year          = {2021},
 archivePrefix = {arXiv},
 title         = {ModelingToolkit: A Composable Graph Transformation System For Equation-Based Modeling}
}

@article{DifferentialEquations.jl-2017,
 author        = {Rackauckas, Christopher and Nie, Qing},
 journal       = {The Journal of Open Research Software},
 title         = {DifferentialEquations.jl – A Performant and Feature-Rich Ecosystem for Solving Differential Equations in Julia},
 number        = {1},
 doi           = {10.5334/jors.151},
 keywords      = {Applied Mathematics},
 year          = {2017},
 url           = {https://app.dimensions.ai/details/publication/pub.1085583166 and http://openresearchsoftware.metajnl.com/articles/10.5334/jors.151/galley/245/download/},
 volume        = {5}
}

@misc{PythonCall.jl,
 author        = {Rowley, Christopher},
 year          = {2022},
 url           = {https://github.com/JuliaPy/PythonCall.jl},
 title         = {PythonCall.jl: Python and Julia in harmony}
}

@article{10.1145/3511528.3511535,
 pages         = {92–96},
 issue_date    = {September 2021},
 author        = {Gowda, Shashi and Ma, Yingbo and Cheli, Alessandro and Gw\'{o}\'{z}zd\'{z}, Maja and Shah, Viral B. and Edelman, Alan and Rackauckas, Christopher},
 month         = {jan},
 journal       = {ACM Commun. Comput. Algebra},
 title         = {High-Performance Symbolic-Numerics via Multiple Dispatch},
 publisher     = {Association for Computing Machinery},
 number        = {3},
 doi           = {10.1145/3511528.3511535},
 address       = {New York, NY, USA},
 year          = {2022},
 url           = {https://doi.org/10.1145/3511528.3511535},
 volume        = {55},
 numpages      = {5},
 abstract      = {As mathematical computing becomes more democratized in high-level languages, high-performance symbolic-numeric systems are necessary for domain scientists and engineers to get the best performance out of their machine without deep knowledge of code optimization. Naturally, users need different term types either to have different algebraic properties for them, or to use efficient data structures. To this end, we developed Symbolics.jl, an extendable symbolic system which uses dynamic multiple dispatch to change behavior depending on the domain needs. In this work we detail an underlying abstract term interface which allows for speed without sacrificing generality. We show that by formalizing a generic API on actions independent of implementation, we can retroactively add optimized data structures to our system without changing the pre-existing term rewriters. We showcase how this can be used to optimize term construction and give a 113x acceleration on general symbolic transformations. Further, we show that such a generic API allows for complementary term-rewriting implementations. Exploiting this feature, we demonstrate the ability to swap between classical term-rewriting simplifiers and e-graph-based term-rewriting simplifiers. We illustrate how this symbolic system improves numerical computing tasks by showcasing an e-graph ruleset which minimizes the number of CPU cycles during expression evaluation, and demonstrate how it simplifies a real-world reaction-network simulation to halve the runtime. Additionally, we show a reaction-diffusion partial differential equation solver which is able to be automatically converted into symbolic expressions via multiple dispatch tracing, which is subsequently accelerated and parallelized to give a 157x simulation speedup. Together, this presents Symbolics.jl as a next-generation symbolic-numeric computing environment geared towards modeling and simulation.}
}

@article{JSSv107i04,
 pages         = {1--32},
 author        = {Bouchet-Valat, Milan and Kamiński, Bogumił},
 journal       = {Journal of Statistical Software},
 title         = {DataFrames.jl: Flexible and Fast Tabular Data in Julia},
 number        = {4},
 doi           = {10.18637/jss.v107.i04},
 year          = {2023},
 url           = {https://www.jstatsoft.org/index.php/jss/article/view/v107i04},
 volume        = {107},
 abstract      = {DataFrames.jl is a package written for and in the Julia language offering flexible and efficient handling of tabular data sets in memory. Thanks to Julia’s unique strengths, it provides an appealing set of features: Rich support for standard data processing tasks and excellent flexibility and efficiency for more advanced and non-standard operations. We present the fundamental design of the package and how it compares with implementations of data frames in other languages, its main features, performance, and possible extensions. We conclude with a practical illustration of typical data processing operations.}
}

@misc{Roots.jl,
 howpublished  = {\url{https://github.com/JuliaMath/Roots.jl}},
 author        = {Verzani, John},
 year          = {2020},
 title         = {{Roots.jl}: Root finding functions for Julia}
}

@article{Julia-2017,
 pages         = {65--98},
 author        = {Bezanson, Jeff and Edelman, Alan and Karpinski, Stefan and Shah, Viral B},
 journal       = {SIAM {R}eview},
 title         = {Julia: A fresh approach to numerical computing},
 publisher     = {SIAM},
 number        = {1},
 doi           = {10.1137/141000671},
 year          = {2017},
 url           = {https://epubs.siam.org/doi/10.1137/141000671},
 volume        = {59}
}

@article{Attractors.jl,
 author        = {Datseris, George and Rossi, Kalel  Luiz and Wagemakers, Alexandre},
 month         = {July},
 journal       = {Chaos: An Interdisciplinary Journal of Nonlinear Science},
 title         = {Framework for global stability analysis of dynamical systems},
 publisher     = {{AIP} Publishing},
 number        = {7},
 doi           = {10.1063/5.0159675},
 year          = {2023},
 url           = {https://doi.org/10.1063/5.0159675},
 volume        = {33}
}

@misc{ComplexityMeasures.jl,
 primaryClass  = {id='cs.SE' full_name='Software Engineering' is_active=True alt_name=None in_archive='cs' is_general=False description='Covers design tools, software metrics, testing and debugging, programming environments, etc. Roughly includes material in all of ACM Subject Classes D.2, except that D.2.4 (program verification) should probably have Logics in Computer Science as the primary subject area.'},
 author        = {Datseris, George and Haaga, Kristian  Agasøster},
 eprint        = {2406.05011},
 year          = {2024},
 archivePrefix = {arXiv},
 title         = {ComplexityMeasures.jl: scalable software to unify and accelerate entropy and complexity timeseries analysis}
}
